\documentclass[11pt,aspectratio=169]{beamer}

\usetheme{default}
\usefonttheme{serif}

\setbeamertemplate{navigation symbols}{}
\setbeamertemplate{footline}[frame number]
\setbeamertemplate{blocks}[rounded][shadow=false]
\setbeamertemplate{frametitle continuation}{}
\setbeamertemplate{itemize items}[triangle]

% Notas do palestrante: descomentar uma das linhas para gerar a versao com notas.
% \setbeameroption{show notes}
% \setbeameroption{show notes on second screen=right}

\usepackage[utf8]{inputenc}
\usepackage[T1]{fontenc}
\usepackage{lmodern}
\usepackage[brazilian]{babel}

\usepackage{amsmath,amssymb}

\usepackage{booktabs}
\usepackage{tabularx}
\usepackage{array}
\usepackage{colortbl}
\usepackage{diagbox}

\usepackage{pgfplots}
\usepackage{pgfgantt}
\pgfplotsset{compat=1.18}
\usetikzlibrary{shapes.geometric,arrows.meta,positioning,calc,fit,backgrounds}

\usepackage{appendixnumberbeamer}

\usepackage[alf]{abntex2cite}
\renewcommand{\refname}{Referências}

% Paleta: laranja Bitcoin como primaria, cinzas como secundarias.
% Elemento grafico e cinza por padrao; vira laranja quando e o objeto da assercao.
\definecolor{laranja}{HTML}{F7931A}
\definecolor{laranjaEsc}{HTML}{B36308}
\definecolor{laranjaCla}{HTML}{FBC77E}
\definecolor{cinzaEsc}{HTML}{3A3A3A}
\definecolor{cinza}{HTML}{7A7A7A}
\definecolor{cinzaCla}{HTML}{D5D5D5}

% Excecao a paleta: quando a cor do balao e o proprio dado, ela e representada
% pela cor real. Uso restrito a dimensao cor do produto.
\definecolor{corBranco}{HTML}{F3EADA}
\definecolor{corAmarelo}{HTML}{F2C230}
\definecolor{corLaranja}{HTML}{E8781F}
\definecolor{corVermelho}{HTML}{C62E1F}
\definecolor{corVinho}{HTML}{7E1F16}
\definecolor{corAzul}{HTML}{2E5FA3}
\definecolor{corVerde}{HTML}{3F8F4A}
\definecolor{corRosa}{HTML}{E58BA9}
\definecolor{corPreto}{HTML}{1F1F1F}

\setbeamercolor{structure}{fg=laranja}
\setbeamercolor{alerted text}{fg=laranjaEsc}
\setbeamercolor{block title}{bg=cinzaCla!45,fg=cinzaEsc}
\setbeamercolor{block body}{bg=cinzaCla!18,fg=black}
\setbeamercolor{frametitle}{fg=cinzaEsc}
\setbeamerfont{frametitle}{series=\bfseries}
\setbeamercolor{footline}{fg=cinza}
\setbeamerfont{footline}{size=\scriptsize}

% Texto nunca recebe a cor primaria; ela fica nos elementos.
\setbeamercolor{title}{fg=cinzaEsc}
\setbeamercolor{subtitle}{fg=cinza}
\setbeamercolor{author}{fg=cinzaEsc}
\setbeamercolor{institute}{fg=cinza}
\setbeamercolor{date}{fg=cinza}
\setbeamercolor{section in toc}{fg=cinzaEsc}
\setbeamercolor{section in toc shaded}{fg=cinza!70}
\setbeamertemplate{section in toc shaded}[default][100]

% Filete laranja sob o titulo do frame: a identidade vive num elemento grafico.
\setbeamertemplate{frametitle}{%
  \vskip0.7em%
  \begin{beamercolorbox}[wd=\paperwidth,leftskip=0.9cm,rightskip=0.9cm]{frametitle}%
    \usebeamerfont{frametitle}\insertframetitle\par%
    \vskip2pt%
    {\color{laranja}\rule{\dimexpr\paperwidth-1.8cm\relax}{1.2pt}}%
  \end{beamercolorbox}%
}

\setbeamertemplate{title page}{%
  \vfill
  \begin{center}
    {\usebeamerfont{title}\usebeamercolor[fg]{title}\inserttitle\par}
    \vskip4pt
    {\color{laranja}\rule{0.45\paperwidth}{1.2pt}}\par
    \vskip6pt
    {\usebeamerfont{subtitle}\usebeamercolor[fg]{subtitle}\insertsubtitle\par}
    \vskip2.2em
    {\usebeamerfont{author}\usebeamercolor[fg]{author}\insertauthor\par}
    \vskip0.6em
    {\usebeamerfont{institute}\usebeamercolor[fg]{institute}\insertinstitute\par}
    \vskip0.6em
    {\usebeamerfont{date}\usebeamercolor[fg]{date}\insertdate\par}
  \end{center}
  \vfill
}

\setlength{\abovedisplayskip}{5pt}
\setlength{\belowdisplayskip}{5pt}
\setlength{\abovedisplayshortskip}{1pt}
\setlength{\belowdisplayshortskip}{1pt}

\newcolumntype{L}[1]{>{\raggedright\arraybackslash}p{#1}}
\newcolumntype{R}{>{\raggedright\arraybackslash}X}

% Faixa localizadora da cadeia herdada: #1 e o indice da familia em destaque.
\newcommand{\faixaliteratura}[2]{%
  \def\focoum{elo}\def\focodois{elo}\def\focotres{elo}\def\focoquatro{elo}%
  \ifnum#1<2 \ifnum#2>0 \def\focoum{foco}\fi\fi
  \ifnum#1<3 \ifnum#2>1 \def\focodois{foco}\fi\fi
  \ifnum#1<4 \ifnum#2>2 \def\focotres{foco}\fi\fi
  \ifnum#1<5 \ifnum#2>3 \def\focoquatro{foco}\fi\fi
  \begin{tikzpicture}[font=\footnotesize,
      elo/.style={rectangle,rounded corners=2pt,draw=cinzaCla,thick,fill=white,text=cinza,
                  text width=2.5cm,minimum height=0.62cm,text centered,inner sep=3pt},
      foco/.style={rectangle,rounded corners=2pt,draw=laranja,very thick,fill=laranja!12,
                   text width=2.5cm,minimum height=0.62cm,text centered,inner sep=3pt},
      elop/.style={draw=cinzaCla,-{Stealth[length=2mm]},thick}]

    \node[\focoum] (ww) at (0,0) {Wagner--Whitin};
    \node[\focodois,right=0.55cm of ww] (clsp) {CLSP};
    \node[\focotres,right=0.55cm of clsp] (small) {GLSP / CLSD};
    \node[\focoquatro,right=0.55cm of small] (carry) {\textit{carryover}};

    \path[elop] (ww) -- (clsp);
    \path[elop] (clsp) -- (small);
    \path[elop] (small) -- (carry);
  \end{tikzpicture}%
}

% Par simbolo-significado indivisivel, para o rodape nunca quebrar no meio de uma entrada.
\newcommand{\sy}[2]{\mbox{$#1$~#2}}

% Rodape de notacao: simbolo -> significado, so nos slides de modelagem.
\newcommand{\notacao}[1]{%
  \vfill
  \noindent{\color{cinzaCla}\rule{\textwidth}{0.4pt}}%
  \par\vspace{0.12em}
  {\scriptsize\color{cinza}\linespread{0.84}\selectfont #1\par}%
}

% Linha de descricao das restricoes, imediatamente acima do bloco algebrico.
\newcommand{\restricoes}[1]{%
  {\footnotesize\color{cinzaEsc}#1\par}%
}

\AtBeginSection[]{%
  \begin{frame}{Roteiro}
    \tableofcontents[currentsection,hideallsubsections]
  \end{frame}
}

\title{Dimensionamento e Sequenciamento de Lotes de Produção}
\subtitle{Uma indústria de balões de látex}
\date{Setembro de 2026}
\author{Felipe Capalbo}
\institute{\small Orientador: Victor C. B. Camargo \\[0.5em] Universidade Federal de São Carlos}

\begin{document}

\maketitle

\begin{frame}{Roteiro}
  \setbeamertemplate{section in toc}[sections numbered]
  \tableofcontents[hideallsubsections]
  \note{Trinta minutos. A apresentação começa pelo problema físico da fábrica, passa pela literatura que
  o enquadra, expõe os dois modelos matemáticos, mostra como eles se tornaram um sistema em
  operação e fecha com o argumento central do trabalho e o que ainda falta.}
\end{frame}

% ===========================================================================
\section{Contexto}
% ===========================================================================

\begin{frame}{O vermelho passa de 34\% a 50\% da carteira}

  \vspace{0.9em}
  \begin{center}
    \begin{tikzpicture}
      \begin{axis}[
          width=11.6cm, height=5.6cm,
          ybar stacked, bar width=0.8cm,
          area legend,
          ylabel={\% da carteira}, ylabel style={font=\scriptsize},
          tick label style={font=\scriptsize},
          xmin=0.4, xmax=7.6,
          xtick={1,2,3,4,5,6,7},
          xticklabels={jan,fev,mar,abr,mai,jun,jul},
          ymin=0, ymax=100, ytick={0,50,100},
          axis lines*=left,
          legend style={at={(0.5,1.04)},anchor=south,draw=none,fill=none,
                        legend columns=7,font=\scriptsize,column sep=5pt},
          clip=false,
        ]

        \addplot[draw=none,fill=corVermelho] coordinates {(1,34.1)(2,34.7)(3,32.7)(4,35.7)(5,42.7)(6,50.4)(7,46.0)};
        \addlegendentry{vermelho}
        \addplot[draw=none,fill=corAzul] coordinates {(1,19.3)(2,18.5)(3,19.6)(4,19.7)(5,19.4)(6,15.4)(7,20.2)};
        \addlegendentry{azul}
        \addplot[draw=none,fill=corBranco] coordinates {(1,16.0)(2,15.2)(3,16.9)(4,16.0)(5,13.1)(6,11.0)(7,18.6)};
        \addlegendentry{branco}
        \addplot[draw=none,fill=corAmarelo] coordinates {(1,12.0)(2,12.6)(3,12.5)(4,12.1)(5,9.5)(6,8.8)(7,12.0)};
        \addlegendentry{amarelo}
        \addplot[draw=none,fill=corVerde] coordinates {(1,9.5)(2,8.3)(3,8.7)(4,7.9)(5,6.7)(6,5.2)(7,3.3)};
        \addlegendentry{verde}
        \addplot[draw=none,fill=corRosa] coordinates {(1,4.6)(2,5.6)(3,4.5)(4,4.5)(5,4.2)(6,4.4)(7,0.0)};
        \addlegendentry{rosa}
        \addplot[draw=none,fill=corPreto] coordinates {(1,4.6)(2,5.1)(3,5.0)(4,4.1)(5,4.4)(6,4.7)(7,0.0)};
        \addlegendentry{preto}

        \coordinate (rotJan) at (axis cs:1,17);
        \coordinate (rotJun) at (axis cs:6,25);

      \end{axis}

      \node[font=\scriptsize,text=white] at (rotJan) {34\%};
      \node[font=\scriptsize,text=white] at (rotJun) {50\%};
    \end{tikzpicture}
  \end{center}

  \note{O eixo de meses é o do slide anterior, agora decomposto por cor. Cada coluna é a
  participação de cada cor na carteira, por mês de vencimento, e todas somam cem por cento: o
  total foi normalizado justamente para que só a composição apareça.

  O vermelho sai de trinta e quatro por cento em janeiro e chega a cinquenta em junho, dezoito
  pontos percentuais entre o menor e o maior mês. O verde faz o caminho inverso, de nove e meio
  para três. Volume igual não significa, portanto, conjunto igual de SKU.

  Daí saem duas consequências. A decisão de capacidade trabalha com o agregado em quilogramas e
  não informa qual cor entra em qual máquina. E a matriz de tempos de troca de cor é percorrida de
  um jeito diferente a cada mês, porque a sequência que a carteira pede muda com ela. São dois
  problemas distintos sobre o mesmo parque, e é por isso que a decisão de quanto produzir e a
  decisão de em que ordem produzir são tomadas em níveis diferentes.

  Uma observação sobre o dado: a carteira rareia à medida que o vencimento se afasta da data de
  corte, e em julho as cores de menor participação ainda não têm pedido colocado. O que a última
  coluna representa é a parcela vendida com mais antecedência.}
\end{frame}

% ---------------------------------------------------------------------------
\begin{frame}{A carga semanal varia de 37\% a 135\% da capacidade}

  \vspace{0.3em}
  \begin{center}
    \begin{tikzpicture}
      \begin{axis}[
          width=12.4cm, height=6.0cm,
          ybar stacked, bar width=0.128cm,
          ylabel={\% da capacidade}, ylabel style={font=\scriptsize},
          tick label style={font=\scriptsize},
          xmin=0.1, xmax=70.9,
          xtick={1,14,27,41,54,67},
          xticklabels={jan/24,abr,jul,out,jan/25,abr},
          ymin=0, ymax=150, ytick={0,50,100,150},
          ymajorgrids, grid style={cinzaCla!60},
          axis lines*=left,
          clip=false,
        ]

        \addplot[draw=none,fill=cinza!60] coordinates {
          (1,100.0) (2,53.5) (3,53.3) (4,52.9) (5,57.1) (6,60.8) (7,70.8) (8,70.8) (9,63.8) (10,69.6) (11,69.6) (12,69.9) (13,69.6) (14,71.3)
          (15,71.0) (16,71.7) (17,71.6) (18,73.9) (19,73.3) (20,73.6) (21,75.1) (22,75.9) (23,87.4) (24,87.5) (25,86.6) (26,85.9) (27,62.2) (28,61.2)
          (29,60.7) (30,61.2) (31,68.5) (32,75.4) (33,74.7) (34,74.2) (35,74.2) (36,82.7) (37,83.0) (38,82.3) (39,82.0) (40,83.7) (41,84.0) (42,84.8)
          (43,85.7) (44,83.6) (45,81.7) (46,82.2) (47,81.3) (48,81.2) (49,36.8) (50,36.8) (51,37.0) (52,92.1) (53,100.0) (54,100.0) (55,54.0) (56,56.2)
          (57,58.0) (58,66.2) (59,75.1) (60,75.1) (61,66.4) (62,68.0) (63,70.0) (64,68.4) (65,67.6) (66,73.1) (67,76.0) (68,73.8) (69,74.8) (70,37.6)};

        \addplot[draw=none,fill=cinzaCla!45] coordinates {
          (1,0.0) (2,46.5) (3,46.7) (4,47.1) (5,42.9) (6,39.2) (7,29.2) (8,29.2) (9,36.2) (10,30.4) (11,30.4) (12,30.1) (13,30.4) (14,28.7)
          (15,29.0) (16,28.3) (17,28.4) (18,26.1) (19,26.7) (20,26.4) (21,24.9) (22,24.1) (23,12.6) (24,12.5) (25,13.4) (26,14.1) (27,37.8) (28,38.8)
          (29,39.3) (30,38.8) (31,31.5) (32,24.6) (33,25.3) (34,25.8) (35,25.8) (36,17.3) (37,17.0) (38,17.7) (39,18.0) (40,16.3) (41,16.0) (42,15.2)
          (43,14.3) (44,16.4) (45,18.3) (46,17.8) (47,18.7) (48,18.8) (49,63.2) (50,63.2) (51,63.0) (52,7.9) (53,0.0) (54,0.0) (55,46.0) (56,43.8)
          (57,42.0) (58,33.8) (59,24.9) (60,24.9) (61,33.6) (62,32.0) (63,30.0) (64,31.6) (65,32.4) (66,26.9) (67,24.0) (68,26.2) (69,25.2) (70,62.4)};

        \addplot[draw=none,fill=laranja] coordinates {
          (1,32.2) (2,0.0) (3,0.0) (4,0.0) (5,0.0) (6,0.0) (7,0.0) (8,0.0) (9,0.0) (10,0.0) (11,0.0) (12,0.0) (13,0.0) (14,0.0)
          (15,0.0) (16,0.0) (17,0.0) (18,0.0) (19,0.0) (20,0.0) (21,0.0) (22,0.0) (23,0.0) (24,0.0) (25,0.0) (26,0.0) (27,0.0) (28,0.0)
          (29,0.0) (30,0.0) (31,0.0) (32,0.0) (33,0.0) (34,0.0) (35,0.0) (36,0.0) (37,0.0) (38,0.0) (39,0.0) (40,0.0) (41,0.0) (42,0.0)
          (43,0.0) (44,0.0) (45,0.0) (46,0.0) (47,0.0) (48,0.0) (49,0.0) (50,0.0) (51,0.0) (52,0.0) (53,20.7) (54,35.1) (55,0.0) (56,0.0)
          (57,0.0) (58,0.0) (59,0.0) (60,0.0) (61,0.0) (62,0.0) (63,0.0) (64,0.0) (65,0.0) (66,0.0) (67,0.0) (68,0.0) (69,0.0) (70,0.0)};

        \coordinate (capIni) at (axis cs:0.1,100);
        \coordinate (capFim) at (axis cs:70.9,100);
        \coordinate (rotCap) at (axis cs:57.5,110);
        \coordinate (rotCarga) at (axis cs:2.5,30);
        \coordinate (rotOcio) at (axis cs:2.5,85);
        \coordinate (rotExc) at (axis cs:47,128);
        \coordinate (alvoExc) at (axis cs:52.3,124);

      \end{axis}

      \draw[cinzaEsc,dashed,thick] (capIni) -- (capFim);
      \node[font=\scriptsize,text=cinzaEsc,anchor=west] at (rotCap) {capacidade};
      \node[font=\scriptsize,text=white,anchor=west] at (rotCarga) {carga};
      \node[font=\scriptsize,text=cinzaEsc,anchor=west] at (rotOcio) {ociosidade};
      \node[font=\scriptsize,text=laranjaEsc,anchor=east] at (rotExc) {excedente};
      \draw[laranjaEsc,-{Stealth[length=2mm]}] ($(rotExc)+(2pt,0)$) -- (alvoExc);
    \end{tikzpicture}
  \end{center}

  \note{A série é a razão entre as horas necessárias para atender a demanda da semana e as horas
  disponíveis das vinte e oito máquinas de moldagem, sobre setenta semanas de histórico escalado
  por um fator aleatório para descaracterizar a grandeza. A média é de setenta e três por cento e a
  faixa vai de trinta e sete a cento e trinta e cinco.

  A parte escura é a carga, a parte clara é a capacidade que sobra na semana e a parte laranja é o
  que não cabe nela. Três semanas das setenta passam de cem por cento e quatro passam de noventa;
  as três estão na virada de ano, quando o calendário de turnos encolhe por parada programada e a
  demanda não encolhe junto.

  O parque de moldagem é fixo e o calendário de turnos é praticamente invariante ao longo do ano,
  porque dimensionar a fábrica para o máximo deixaria o investimento ocioso na maior parte do
  horizonte. Atender as semanas de maior carga exige antecipar produção nas semanas de ociosidade
  e carregar o estoque até lá: é assim que a folga sai da capacidade e passa para o estoque, que
  deixa de ser consequência do plano e vira variável de decisão. É também o que amarra a previsão
  de demanda ao dimensionamento de lotes, porque antecipar contra um máximo só é possível se o
  modelo enxerga esse máximo.

  Convém registrar o que a conta não inclui: as horas de \textit{setup} e a compatibilidade entre
  máquina e balão. A utilização efetiva é, portanto, maior que a apurada aqui.}
\end{frame}

% ---------------------------------------------------------------------------
\begin{frame}{A moldagem concentra os dois \textit{setups} do problema}

  \begin{center}
    \resizebox{!}{0.68\textheight}{%
    \begin{tikzpicture}[
        process/.style={rectangle,draw=cinza,thick,fill=white,text centered,
                        minimum height=1.5cm,text width=2.5cm,font=\small},
        inventory/.style={isosceles triangle,shape border rotate=180,draw=cinza,thick,
                          fill=cinza!15,text centered,minimum height=2.0cm,text width=2.0cm,
                          inner ysep=0pt,text depth=1.2em,font=\small},
        critical/.style={process,very thick,draw=laranjaEsc,fill=laranjaEsc!10,font=\small\bfseries},
        colorstep/.style={process,very thick,draw=laranja,fill=laranja!12},
        arrow/.style={draw=cinza,-{Stealth[length=3mm]},thick},
        note/.style={font=\footnotesize\itshape,text width=4.2cm,text centered}
      ]

      \node[process] (receb) {Recebimento};
      \node[inventory,right=1.1cm of receb] (elatex) {Látex};
      \node[process,right=1.1cm of elatex] (prevul) {Pré-vulcanização};
      \node[colorstep,right=1.1cm of prevul] (mistura) {Mistura e pigmentação};

      \node[critical,below=2.6cm of mistura] (moldagem) {Moldagem e vulcanização};
      \node[process,left=1.1cm of moldagem] (talco) {Remoção de talco};
      \node[process,left=1.1cm of talco] (embal) {Embalagem};
      \node[inventory,left=1.1cm of embal] (epa) {Produto acabado};

      \path[arrow] (receb) -- (elatex);
      \path[arrow] (elatex) -- (prevul);
      \path[arrow] (prevul) -- (mistura);
      \path[arrow] (mistura) -- (moldagem);
      \path[arrow] (moldagem) -- (talco);
      \path[arrow] (talco) -- (embal);
      \path[arrow] (embal) -- (epa);

      \node[note,text=laranja,above=0.3cm of mistura] {definição da cor};
      \node[note,text=laranjaEsc,below=0.3cm of moldagem] {objeto do estudo \\ \textit{setup} de formato e de cor};

    \end{tikzpicture}}
  \end{center}

  \note{O látex chega todo dia porque é perecível, e fica em tanques. Na pré-vulcanização ele
  recebe os agentes químicos que dão elasticidade. Na mistura entram pigmento e aditivo de
  acabamento, e é exatamente aí que o balão deixa de ser genérico e vira cor-específico.

  Na moldagem os moldes imergem no composto, passam por talco, vulcanizam em estufa e são
  desmoldados. Os estágios anteriores têm capacidade excedente, então modelar apenas a capacidade
  da moldagem é suficiente. É também na moldagem que ocorrem os dois setups do problema, e é por
  isso que ela é o objeto deste estudo. Os dois setups são o próximo slide.}
\end{frame}

% ---------------------------------------------------------------------------
\begin{frame}{A troca de formato ocupa 60\% do turno; a de cor, 11\%}

  \vspace{0.5em}
  \begin{center}
    \begin{tikzpicture}[x=0.865cm,y=1.42cm,font=\small,
        chamada/.style={align=center,inner sep=1pt,text=cinzaEsc},
        guia/.style={draw=cinza,thin}]

      \fill[cinzaCla!30] (0,0) rectangle (8,1);
      \fill[cinzaCla!30] (8,0) rectangle (16,1);

      \fill[corBranco]   (0,0.13)    rectangle (2.20,0.87);
      \fill[laranja]     (2.20,0)    rectangle (3.05,1);
      \fill[corAmarelo]  (3.05,0.13) rectangle (5.30,0.87);
      \fill[laranja]     (5.30,0)    rectangle (6.15,1);
      \fill[corVermelho] (6.15,0.13) rectangle (8,0.87);

      \fill[laranjaEsc]  (8,0)       rectangle (12.80,1);
      \fill[corVermelho] (12.80,0.13) rectangle (16,0.87);

      \draw[cinza!50] (0,0) rectangle (8,1);
      \draw[cinza!50] (8,0) rectangle (16,1);
      \draw[cinzaEsc,thick] (8,-0.18) -- (8,1.18);

      \draw[guia] (2.625,1) -- (2.625,1.45);
      \node[chamada,anchor=south] at (2.625,1.50)
        {{\footnotesize troca de cor}\\[0.15em]
         {\bfseries 0,85 h}\;{\scriptsize\color{cinza}0,33--3,40}\\[0.05em]
         {\scriptsize\color{cinza}decidida no dia}};

      \draw[guia] (10.40,1) -- (10.40,1.45);
      \node[chamada,anchor=south] at (10.40,1.50)
        {{\footnotesize troca de formato}\\[0.15em]
         {\bfseries 4,8 h}\;{\scriptsize\color{cinza}1,5--6,3}\\[0.05em]
         {\scriptsize\color{cinza}decidida na semana}};

      \draw[cinza,{Bar[width=5pt]}-{Bar[width=5pt]}] (0,-0.45) -- (8,-0.45);
      \draw[cinza!50,{Bar[width=5pt]}-{Bar[width=5pt]}] (8,-0.45) -- (16,-0.45);
      \node[font=\scriptsize,text=cinzaEsc,anchor=north,fill=white,inner sep=2pt] at (4,-0.36) {turno 8 h};
      \node[font=\scriptsize,text=cinza,anchor=north,fill=white,inner sep=2pt] at (12,-0.36) {turno};

    \end{tikzpicture}
  \end{center}

  \note{A escala de tempo dos dois \textit{setups} é a característica que estrutura o trabalho inteiro.
  Na planilha da fábrica, a troca de formato vai de 1,5 a 6,3 horas entre os dezessete balões, com
  mediana de 4,8 horas, e depende do par de formatos e da máquina. A troca de cor vai de 0,33 a
  3,40 horas, com mediana de 0,85 hora. A razão entre as duas medianas é de 5,6 vezes. As faixas
  se sobrepõem, e por isso a comparação apresentada é entre as medianas, não entre os extremos.

  O que torna essa diferença decisiva é o turno de oito horas. A troca de formato consome cerca de
  sessenta por cento da capacidade do turno, de modo que executá-la no meio dele não é operação
  real: ela ocorre na virada, quando a máquina já está parando, e é decidida com semanas de
  antecedência, no plano semanal. A troca de cor cabe entre duas corridas do mesmo balão, consome
  em torno de um décimo do turno e é decidida no dia, dentro do turno corrente.

  A troca de cor é ainda direcional, com transições proibidas; a matriz completa aparece adiante,
  no modelo por turno. O ponto deste slide é apenas a diferença de escala e de cadência, porque é
  dela que vem a decomposição em dois modelos: um semanal, que decide formato e volume por máquina,
  e outro por turno, que sequencia as cores sobre o formato já fixado.}
\end{frame}

% ---------------------------------------------------------------------------
\begin{frame}{A carteira cai de 95\% a 20\% da demanda no horizonte}

  \vspace{0.2em}
  \begin{center}
    \begin{tikzpicture}
      \begin{axis}[
          width=11.6cm, height=5.5cm,
          xlabel={semana}, ylabel={mil kg},
          xlabel style={font=\scriptsize}, ylabel style={font=\scriptsize},
          tick label style={font=\scriptsize},
          xtick={1,2,3,4,5,6,7,8,9,10,11,12,13},
          xmin=1, xmax=13,
          ymin=0, ymax=56, ytick={0,20,40},
          ymajorgrids, grid style={cinza!20},
          axis lines*=left,
          clip=false,
        ]

        \addplot[draw=none,fill=cinzaCla] coordinates {
          (1,40.7) (2,40.7) (3,40.7) (4,40.7) (5,43.5) (6,46.2) (7,46.2)
          (8,46.2) (9,48.3) (10,52.5) (11,52.5) (12,52.5) (13,52.5)} \closedcycle;

        \addplot[draw=none,fill=laranja] coordinates {
          (1,38.5) (2,33.9) (3,29.9) (4,26.4) (5,24.8) (6,23.3) (7,20.4)
          (8,17.9) (9,16.5) (10,15.9) (11,14.0) (12,12.2) (13,10.7)} \closedcycle;

        \node[font=\small\bfseries,text=cinzaEsc,anchor=west] at (axis cs:1.35,21.5) {carteira};
        \node[font=\scriptsize,text=cinzaEsc,anchor=west]     at (axis cs:1.35,16.0) {falta $=$ atraso};

        \node[font=\small\bfseries,text=cinzaEsc,anchor=west] at (axis cs:7.1,42.0) {previsão};
        \node[font=\scriptsize,text=cinzaEsc,anchor=west]     at (axis cs:7.1,36.5) {falta $=$ venda perdida};

        \draw[laranjaEsc,thick] (axis cs:1.55,45.5) -- (axis cs:1.08,39.6);
        \node[font=\small\bfseries,text=laranjaEsc,anchor=south west] at (axis cs:1.45,45.5) {95\%};

        \draw[laranjaEsc,thick] (axis cs:12.25,19.0) -- (axis cs:12.85,11.6);
        \node[font=\small\bfseries,text=laranjaEsc,anchor=south east] at (axis cs:12.35,18.6) {20\%};

      \end{axis}
    \end{tikzpicture}
  \end{center}

  \note{A demanda que alimenta o planejamento não chega por uma via única. Parte está confirmada em
  pedidos com quantidade, cor e prazo acordados, tratada aqui como carteira; o restante é previsão.
  A carteira apurada na fábrica reúne oitocentos e sessenta e três pedidos, dois mil e oitenta e
  dois itens e vinte clientes, e os valores do gráfico estão escalados por um fator aleatório.

  A proporção entre as duas correntes varia ao longo do horizonte. Na primeira semana noventa e
  cinco por cento da demanda já está vendida, e na décima terceira restam vinte por cento: à
  medida que o horizonte se afasta a carteira rareia e a previsão predomina. A área total ainda
  cresce, porque o horizonte entra na subida sazonal, de modo que o que encolhe é a participação
  da carteira, não o seu volume apenas.

  A distinção não é nominal. A falta contra carteira é atraso, transportado adiante e ainda
  devido; a falta contra previsão é venda perdida, que não se recupera. A operação é híbrida entre
  produção contra pedido e produção para estoque sobre o mesmo parque de máquinas, com as duas
  correntes disputando a mesma capacidade, e o modelo semanal carrega duas variáveis de falta
  distintas exatamente por causa dessa assimetria.}
\end{frame}

% ---------------------------------------------------------------------------

% ---------------------------------------------------------------------------
\begin{frame}{A cadência da decisão separa os dois modelos}

  \vspace{0.6em}
  \begin{center}
    \begin{tikzpicture}[font=\small]

      \foreach \sem/\part in {1/95,2/83,3/73,4/65,5/57,6/50,7/44,8/39,9/34,
                              10/30,11/27,12/23,13/20} {
        \pgfmathsetmacro{\xesq}{(\sem-1)*0.72}
        \pgfmathsetmacro{\xdir}{\xesq+0.66}
        \fill[laranja!\part] (\xesq,1.62) rectangle (\xdir,2.28);
        \node[font=\scriptsize,text=cinza] at ({(\xesq+\xdir)/2},2.50) {\sem};
      }

      \node[font=\small\bfseries,text=laranjaEsc,anchor=west] at (0,2.92) {carteira 95\%};
      \node[font=\small\bfseries,text=laranjaEsc,anchor=east] at (9.36,2.92) {20\%};

      \fill[cinzaCla!45] (0,1.62) -- (0.66,1.62) -- (9.36,0.74) -- (0,0.74) -- cycle;
      \draw[laranjaEsc,thick] (0,1.62) rectangle (0.66,2.28);

      \foreach \dia in {0,1,2,3,4,5} {
        \foreach \turno in {0,1,2} {
          \pgfmathsetmacro{\xesq}{\dia*1.56+\turno*0.46}
          \pgfmathsetmacro{\xdir}{\xesq+0.40}
          \fill[laranjaEsc] (\xesq,0.08) rectangle (\xdir,0.68);
        }
      }

      \node[anchor=east,font=\small\bfseries]                      at (-0.22,2.10) {Semanal};
      \node[anchor=east,font=\scriptsize,text=cinza,
            text width=1.9cm,align=right]                          at (-0.22,1.70) {quantidade por máquina};
      \node[anchor=east,font=\small\bfseries]                      at (-0.22,0.56) {Turno};
      \node[anchor=east,font=\scriptsize,text=cinza,
            text width=1.9cm,align=right]                          at (-0.22,0.18) {sequência de cores};

      \node[anchor=west,font=\scriptsize,text=cinza] at (9.55,0.38) {18 turnos};

      \draw[cinza,-{Stealth[length=2.5mm]},thick]
        (9.36,3.55) .. controls (6.5,4.15) and (2.5,4.15) .. (0.05,3.55);
      \node[font=\scriptsize,text=cinza] at (4.68,4.22) {horizonte rolante};

    \end{tikzpicture}
  \end{center}

  \note{As decisões não compartilham a mesma cadência, e é essa diferença que fixa a fronteira
  entre os dois modelos. Quanto produzir de cada balão, em que máquina e em que semana é decisão
  revisada semanalmente, sobre um horizonte que precisa conter a subida sazonal. Em que turno e em
  que ordem as cores entram em cada máquina é decisão da semana corrente, condicionada à carteira
  já existente.

  A intensidade das células retoma o gráfico anterior: a participação da carteira cai de noventa e
  cinco por cento na primeira semana para vinte por cento na décima terceira. É o argumento
  decisivo contra o modelo único, que obrigaria a decidir hoje a cor que a máquina produz em
  dezembro, quando a carteira de pedidos de dezembro ainda não existe. Decidir cedo demais não é
  decidir melhor, é decidir sobre informação que ainda não chegou.

  Só a primeira semana é detalhada em dezoito turnos, três por dia ao longo de seis dias; as
  demais permanecem agregadas até a revisão seguinte, num esquema de horizonte rolante. Cada uma
  das três escalas encaixadas, semana, turno e posição dentro do turno, responde a um fenômeno
  físico do processo e não a conveniência de modelagem: a semana é o intervalo de revisão do
  plano, o turno é a fronteira das trocas de forma e a posição é o que determina qual tempo de
  limpeza entre cores é pago.

  O que desce da semana para o turno é uma única grandeza, o estoque que cada balão deve ter ao
  fim da semana. O que sobe de volta, na revisão seguinte, é o estado real. A decomposição reduz o
  tamanho de cada modelo, mas o tamanho é consequência, não motivo.}
\end{frame}

% ===========================================================================
\section{Literatura}
% ===========================================================================

\begin{frame}{Ao balanço clássico acrescenta-se a venda perdida}

  \begin{center}
    \faixaliteratura{1}{2}\\[0.2em]
    {\scriptsize\color{cinza}\cite{wagnerwhitin1958} \quad \cite{karimi2003}}
  \end{center}

  \begin{center}
    \begin{tikzpicture}[font=\footnotesize,
        cx/.style={rectangle,rounded corners=2pt,draw=cinza,thick,fill=white,inner sep=6pt},
        cn/.style={rectangle,rounded corners=2pt,draw=laranja,very thick,fill=laranja!12,inner sep=6pt}]

      \node[cx] (lit) at (0,0) {$\begin{aligned}
          I_{t-1} + x_{t} &= I_{t} + d_{t} \\
          \textstyle\sum_{i} \bigl(a_{i} x_{it} + \sigma_{i} y_{it}\bigr) &\le b_{t}
        \end{aligned}$};
      \node[font=\footnotesize,text=cinza,above=2pt of lit] {literatura};

      \node[cn,below=0.45cm of lit] (nos) {$\begin{aligned}
          I_{i,t-1} + \textstyle\sum_{j \in J_i} x_{ijt}
            &= I_{it} + \bigl(o_{it} + a_{i,t-1} - a_{it}\bigr) + \bigl(\hat{d}_{it} - l_{it}\bigr) \\
          \textstyle\sum_{i \in I} \bigl( x_{ijt}/p_{ij} + \tau_{i} y_{ijt} \bigr) &\le \kappa_{jt}
        \end{aligned}$};
      \node[font=\footnotesize,text=laranjaEsc,anchor=south east] at ([yshift=2pt]nos.north east)
        {este trabalho};

      \draw[-{Stealth[length=3mm]},thick,cinza] (lit.south) -- (nos.north);

      \node[font=\scriptsize,text=cinzaEsc,anchor=north] at ([yshift=-4pt]nos.south)
        {balde é a semana-máquina \qquad $l_{it}$ escolhe o que não vender};
    \end{tikzpicture}
  \end{center}

  \notacao{%
    \sy{i}{balão} · \sy{j}{máquina} · \sy{t}{semana} · \sy{I}{estoque} · \sy{x}{produção} · 
    \sy{y}{\textit{setup}} · \sy{o}{carteira} · \sy{\hat{d}}{previsão} · \sy{a}{atraso} · 
    \sy{l}{venda perdida} · \sy{p}{taxa} · \sy{\tau}{horas de \textit{setup}} · \sy{\kappa}{horas disponíveis}}

  \note{A cadeia tem quatro elos, e a caixa em destaque no alto indica em que ponto dela a
  apresentação está. Os dois primeiros elos são \textit{big-bucket} e sustentam o nível semanal.

  De Wagner e Whitin herda-se o esqueleto: o balanço recursivo de estoque, período a período, e a
  ideia de cobrar o \textit{setup} como carga fixa disparada por uma binária. Do CLSP herda-se a
  capacidade por período em unidade de tempo, e o fato de o \textit{setup} consumir essa
  capacidade em vez de aparecer só como custo.

  Duas coisas mudam no balanço, e as duas vêm do problema real. O custo de estoque sai da função
  objetivo, porque a empresa não tem armazenagem relevante e o custo que importa é a oportunidade
  perdida. E a demanda deixa de ser sempre atendida: a capacidade pode não dar, e o modelo precisa
  escolher o que deixar de vender. É o $l$ de venda perdida entrando no balanço, ao lado do
  atraso de carteira, que é transportado de uma semana para a seguinte.

  Na capacidade o que muda é o balde. No CLSP ele é o período inteiro; aqui é o par
  semana-máquina, e o $p$ traz índice de máquina e de balão. É isso que caracteriza máquinas
  paralelas não relacionadas: cada par tem sua própria taxa, e não existe máquina uniformemente
  mais rápida que outra.

  O que este nível ainda não faz é ordenar dentro do período. Para o formato isso basta, porque a
  troca de formato consome o mesmo tempo venha de onde vier. Para a cor, não basta.}
\end{frame}

% ---------------------------------------------------------------------------
\begin{frame}{À sequência do GLSP acrescenta-se a proibição}

  \begin{center}
    \faixaliteratura{3}{4}\\[0.2em]
    {\scriptsize\color{cinza}\cite{fleischmannmeyr1997} \quad \cite{haase1996} \quad \cite{suriestadtler2003}}
  \end{center}

  \vspace{0.3em}
  \begin{center}
    \begin{tikzpicture}[font=\small,
        cx/.style={rectangle,rounded corners=2pt,draw=cinza,thick,fill=white,inner sep=7pt},
        cn/.style={rectangle,rounded corners=2pt,draw=laranja,very thick,fill=laranja!12,inner sep=7pt}]

      \node[cx] (lit) at (0,0) {$\begin{aligned}
          \textstyle\sum_{i} y_{is} &= 1 \\
          z_{ii's} &\ge y_{i,s-1} + y_{i's} - 1 \\
          \delta_{it} &\ge y_{it} - y_{i,t-1}
        \end{aligned}$};
      \node[font=\footnotesize,text=cinza,above=2pt of lit] {literatura};

      \node[cn,right=1.0cm of lit] (nos) {$\begin{aligned}
          \textstyle\sum_{k \in K_i} \nu_{kjsn} &= u_{js} \\
          \zeta_{kk'jsn} &\ge \nu_{kjs,n-1} + \nu_{k'jsn} - 1 \\
          \nu_{kjs,n-1} + \nu_{k'jsn} &\le 1 \quad {\scriptstyle (k,k') \notin A} \\
          \bigl| \varphi_{kjs} - \varphi_{kj,s-1} \bigr| &\le u_{js}
        \end{aligned}$};
      \node[font=\footnotesize,text=laranjaEsc,above=2pt of nos] {este trabalho};

      \draw[-{Stealth[length=3mm]},thick,cinza] (lit) -- (nos);

      \node[font=\scriptsize,text=cinza,anchor=north,text width=4.6cm,align=center]
        at ([yshift=-5pt]lit.south) {microperíodo ordena \\ \textit{setup} só na mudança};
      \node[font=\scriptsize,text=cinzaEsc,anchor=north,text width=5.4cm,align=center]
        at ([yshift=-5pt]nos.south) {microperíodo é a posição no turno \\ par fora de $A$ é inviável};
    \end{tikzpicture}
  \end{center}

  \notacao{%
    \sy{j}{máquina} · \sy{s}{turno} · \sy{n}{posição} · \sy{k}{cor} · \sy{K_i}{cores do balão} · 
    \sy{A}{transições permitidas} · \sy{\nu}{cor na posição} · \sy{u}{turno ativo} · 
    \sy{\zeta}{troca de cor} · \sy{\varphi}{cor ao fim do turno}}

  \note{Os dois últimos elos são \textit{small-bucket} e sustentam o nível operacional.

  Quando o \textit{setup} depende da sequência é preciso saber a ordem, e a literatura resolve
  isso descendo para microperíodos. O GLSP e o CLSD fornecem duas restrições: um único estado por
  microperíodo, e a detecção da transição pelo produto de dois estados consecutivos, linearizado
  naquela desigualdade com o menos um. O \textit{carryover} fornece a terceira: o \textit{setup}
  deixa de ser cobrado em todo período ativo e passa a ser pago só quando o estado muda.

  As três são herdadas com tradução direta. O microperíodo é a posição dentro do turno, o item é a
  cor, e a forma já está fixada para aquele turno. O número de posições é pequeno porque o lote
  mínimo e o tempo de limpeza limitam quantas corridas de cor cabem num turno de oito horas.

  A terceira linha da direita é o que a literatura não fornece. Quando o par não pertence ao
  conjunto de transições permitidas, os dois estados não podem ser consecutivos. Isso normalmente
  é modelado com custo alto; aqui é restrição, porque depois de preto voltar para branco não é
  questão de tempo, é inviável sem purgar o tanque. Essa linha carrega o \textit{setup}
  direcional, que é a característica sem equivalente nas quatro classes revisadas.

  A quarta linha é o \textit{carryover} na forma que este trabalho usa: o $\varphi$ carrega a cor
  montada de um turno para o seguinte, e quando a máquina não produz o $u$ é zero e a cor
  permanece. Sem esse mecanismo, descer para o turno multiplicaria o custo de \textit{setup} pelo
  número de turnos, e uma máquina parada atravessaria a troca sem pagá-la.}
\end{frame}

% ---------------------------------------------------------------------------
\begin{frame}{Três trabalhos análogos resolvem parte do problema}

  \vspace{0.6em}
  \begin{columns}[T,onlytextwidth]
    \begin{column}{0.31\textwidth}
      \begin{center}
        \begin{tikzpicture}[font=\footnotesize]
          \node[rectangle,draw=cinza,thick,fill=cinza!10,minimum width=2.2cm,
                minimum height=0.8cm] (a) at (0,1.0) {xarope};
          \node[rectangle,draw=laranjaEsc,very thick,fill=laranjaEsc!10,minimum width=2.2cm,
                minimum height=0.8cm] (b) at (0,0) {envase};
          \draw[-{Stealth[length=2mm]},thick,cinza] (a) -- (b);
        \end{tikzpicture}
      \end{center}
      \vspace{0.4em}
      \begin{center}\footnotesize\textbf{Bebidas} \\ {\scriptsize\cite{ferreira2010}}\end{center}
    \end{column}
    \begin{column}{0.31\textwidth}
      \begin{center}
        \begin{tikzpicture}[font=\footnotesize]
          \node[rectangle,draw=cinza,thick,fill=cinza!10,minimum width=2.2cm,
                minimum height=0.8cm] (a) at (0,1.0) {mistura};
          \node[rectangle,draw=laranjaEsc,very thick,fill=laranjaEsc!10,minimum width=0.6cm,
                minimum height=0.8cm] (b1) at (-0.8,0) {};
          \node[rectangle,draw=laranjaEsc,very thick,fill=laranjaEsc!10,minimum width=0.6cm,
                minimum height=0.8cm] (b2) at (0,0) {};
          \node[rectangle,draw=laranjaEsc,very thick,fill=laranjaEsc!10,minimum width=0.6cm,
                minimum height=0.8cm] (b3) at (0.8,0) {};
          \draw[-{Stealth[length=2mm]},thick,cinza] (a) -- (b1);
          \draw[-{Stealth[length=2mm]},thick,cinza] (a) -- (b2);
          \draw[-{Stealth[length=2mm]},thick,cinza] (a) -- (b3);
        \end{tikzpicture}
      \end{center}
      \vspace{0.4em}
      \begin{center}\footnotesize\textbf{Fiação} \\ {\scriptsize\cite{camargo2014}}\end{center}
    \end{column}
    \begin{column}{0.31\textwidth}
      \begin{center}
        \begin{tikzpicture}[font=\footnotesize]
          \foreach \i/\o in {0/15,1/40,2/65,3/90}
            \node[rectangle,draw=laranja,thick,fill=laranja!\o,
                  minimum width=0.5cm,minimum height=0.8cm] at (-1.05+\i*0.7,0.5) {};
          \draw[-{Stealth[length=2mm]},thick,cinza] (-1.5,-0.35) -- (1.5,-0.35);
          \node[font=\scriptsize,text=cinza] at (0,-0.85) {claro $\to$ escuro};
        \end{tikzpicture}
      \end{center}
      \vspace{0.4em}
      \begin{center}\footnotesize\textbf{\textit{Block planning}} \\ {\scriptsize\cite{gunther2006}}\end{center}
    \end{column}
  \end{columns}

  \note{Três analogias, todas parciais.

  Bebidas é a mais estudada: xarope preparado num tanque, envasado depois em sabores e tamanhos.
  Composto pigmentado é o xarope, moldagem é o envase. Ferreira, Morabito e Rangel modelam
  exatamente isso com setup dependente de sequência.

  Fiação acrescenta o que bebidas não tem: uma mistura de fibras alimenta várias máquinas em
  paralelo, e todas as ativas num microperíodo precisam usar a mesma mistura. É o gabarito formal
  caso a pigmentação venha a ser modelada como estágio próprio.

  Block planning é a regra de sequenciamento: produzir do tom mais claro para o mais escuro para
  evitar limpeza. É o que a fábrica já pratica.

  Duas diferenças mantêm este trabalho fora dessas formulações: elas tratam a troca de mistura
  como setup nulo, enquanto aqui a troca de formato é o setup dominante; e elas usam backorder com
  custo de estoque explícito, enquanto aqui se usa venda perdida e custo de oportunidade puro.}
\end{frame}

% ===========================================================================
\section{Modelos}
% ===========================================================================

\begin{frame}{A sazonalidade só é nítida nos níveis agregados}

  \vspace{0.5em}
  \begin{center}
    \resizebox{\textwidth}{!}{%
    \begin{tikzpicture}[font=\small,
        nd/.style={rectangle,rounded corners=2pt,thick,minimum height=0.75cm},
        n0/.style={nd,draw=laranja,very thick,fill=laranja!20},
        n1/.style={nd,draw=cinza,fill=cinza!14},
        n2/.style={nd,draw=cinza,fill=white},
        ar/.style={draw=cinza,-{Stealth[length=2.5mm]},thick}]

      \node[n0,minimum width=7.0cm] (g) at (0,3.4) {global};
      \node[n1,minimum width=2.0cm] (f1) at (-2.4,1.7) {$f_1$};
      \node[n1,minimum width=2.0cm] (f2) at (0,1.7) {$f_2$};
      \node[n1,minimum width=2.0cm] (f3) at (2.4,1.7) {$f_3$};
      \node[n2,minimum width=1.4cm] (i11) at (-3.1,0) {$i_{11}$};
      \node[n2,minimum width=1.4cm] (i12) at (-1.6,0) {$i_{12}$};
      \node[n2,minimum width=1.4cm] (i21) at (0,0) {$i_{21}$};
      \node[n2,minimum width=1.4cm] (i31) at (2.4,0) {$i_{31}$};

      \path[ar] (g) -- (f1); \path[ar] (g) -- (f2); \path[ar] (g) -- (f3);
      \path[ar] (f1) -- (i11); \path[ar] (f1) -- (i12);
      \path[ar] (f2) -- (i21); \path[ar] (f3) -- (i31);

      \node[anchor=west,font=\footnotesize,text=cinzaEsc] at (4.0,3.4) {sinal forte};
      \node[anchor=west,font=\footnotesize,text=cinzaEsc] at (4.0,1.7) {\textit{prior}};
      \node[anchor=west,font=\footnotesize,text=cinzaEsc] at (4.0,0)  {ruído};

      \draw[-{Stealth[length=3mm]},very thick,cinza] (6.6,3.1) -- (6.6,0.3);
      \node[anchor=west,font=\footnotesize\itshape,text=cinzaEsc,align=left] at (6.9,1.7)
        {$\lambda$ penaliza \\ o desvio};
    \end{tikzpicture}}
  \end{center}

  \note{O problema da previsão aqui é de escala. No nível do balão individual, o histórico é curto e
  ruidoso, e a sazonalidade se perde. No agregado, ela é limpa.

  A ideia é estimar onde o sinal é forte e repassar para baixo como informação a priori. Cada
  nível é estimado com os coeficientes do nível acima funcionando como âncora, e um parâmetro
  lambda penaliza o quanto o filho pode se afastar do pai.

  Uma escolha de modelagem que merece explicação: o modelo é log-linear, multiplicativo, não
  aditivo. No máximo de fim de ano a demanda agregada pode triplicar. Cascatear um incremento
  absoluto de dez mil quilos do nível global para um produto que vende trezentos faria o número
  perder qualquer sentido. Em log, o efeito vira um fator percentual, que desce coerente entre
  escalas muito diferentes — e ainda impede previsão negativa por construção.}
\end{frame}

% ---------------------------------------------------------------------------
\begin{frame}{O nível agregado ancora a estimação do nível abaixo}

  \restricoes{%
    (A) estimação global, sem \textit{prior} \hfill
    (B) estimação ancorada no nível acima \hfill
    (C) previsão do mês alvo}

  \vspace{0.2em}
  {\footnotesize
  \begin{equation*}
    \hat{\beta}^{(0)} = \arg\min_{\beta} \sum_{t \in T}
      \Bigl( \ln(Y_t) - \sum_{k \in K} \beta_{k} X_{kt} \Bigr)^2 \tag{A}
  \end{equation*}

  \begin{equation*}
    \hat{\beta}^{(\ell)} = \arg\min_{\beta}
      \Bigl( \sum_{t \in T} \bigl( \ln(Y^{(\ell)}_{t}) - \textstyle\sum_{k} \beta_{k} X_{kt} \bigr)^2
      + \alert{\lambda \sum_{k \in K} \bigl( \beta_{k} - \hat{\beta}^{(\ell-1)}_{k} \bigr)^2} \Bigr)
      \quad \ell \in \{1,\,2\} \tag{B}
  \end{equation*}

  \begin{equation*}
    \hat{d}_{i,t+h} = B_{i} \cdot \exp\Bigl( \sum_{k \in K} \hat{\beta}^{(2)}_{ik} X_{k,t+h} \Bigr) \tag{C}
  \end{equation*}}

  \vspace{0.3em}
  \begin{center}
    \begin{tikzpicture}[font=\footnotesize]
      \node[rectangle,rounded corners=2pt,draw=laranja,very thick,fill=laranja!12,
            minimum width=2.4cm,minimum height=0.7cm] (d) at (0,0) {$\hat{d}_{i,t+h}$};
      \node[rectangle,rounded corners=2pt,draw=cinza,very thick,fill=cinza!12,
            minimum width=2.6cm,minimum height=0.8cm] (t) at (5.2,0) {semanal};
      \draw[-{Stealth[length=3mm]},very thick,cinza] (d) -- node[above,font=\scriptsize,text=cinza] {$d_{it}$} (t);
    \end{tikzpicture}
  \end{center}

  \notacao{%
    \sy{t}{mês} · \sy{k}{mês do ano} · \sy{f}{formato} · \sy{i}{balão} · \sy{\ell}{nível} · 
    \sy{Y_t}{demanda global} · \sy{Y_{ft}}{demanda do formato} · \sy{y_{it}}{demanda do balão} · 
    \sy{X_{kt}}{\textit{feature} sazonal} · \sy{\beta}{coeficiente sazonal} · 
    \sy{\lambda}{penalização do desvio} · \sy{B_i}{demanda base} · \sy{\hat{d}}{previsão}}

  \note{Três estimações em cascata. A primeira é global, sem \textit{prior}: mínimos quadrados puro
  sobre o log da demanda total. As duas seguintes seguem a mesma forma, com o nível acima
  funcionando como âncora — formato ancorado no global, balão ancorado no formato —, e trazem o
  termo destacado: além de ajustar o próprio histórico, cada nível paga uma penalidade por se
  afastar do coeficiente do pai.

  O comportamento nos extremos ajuda a entender o papel de lambda. Com lambda tendendo ao
  infinito, o balão herda integralmente a sazonalidade do formato e ignora o próprio histórico.
  Com lambda zero, cada balão é estimado sozinho, sobre a série curta e ruidosa — que é
  exatamente o que se quer evitar. A estimação usa mínimos quadrados recursivos, o que permite
  atualizar os coeficientes a cada nova observação sem reestimar tudo.

  A última equação fecha a volta à escala original. O exponencial dos coeficientes granulares é o
  efeito sazonal consolidado daquele balão naquele mês; dividindo a demanda observada por ele
  sobra a demanda base, o que o produto vende fora de sazonalidade, e a previsão futura é a base
  multiplicada pelo efeito esperado do mês alvo.

  O que importa reter é o diagrama: esse d chapéu é literalmente o parâmetro de demanda que entra
  no modelo semanal. E existe um segundo canal entre os dois, que é o alfa — quantos períodos
  futuros o estoque deve cobrir. É o alfa que converte acurácia de previsão em política de
  estoque: quanto melhor a previsão, menor o alfa necessário, maior o giro. Essa é a integração
  que a literatura de \textit{lot sizing} revisada não tem, porque ela recebe a previsão pronta
  de fora.}
\end{frame}

% ---------------------------------------------------------------------------
\begin{frame}{As duas demandas competem pela mesma capacidade}

  \begin{columns}[T,onlytextwidth]
    \begin{column}{0.24\textwidth}
      {\footnotesize\color{cinza}\textbf{Conjuntos}}
      \vspace{0.3em}
      {\scriptsize
      \begin{tabularx}{\linewidth}{@{}L{1.0cm}R@{}}
        $T$   & Semanas. \\
        $I$   & Balões. \\
        $J$   & Máquinas. \\
        $J_i$ & Compatíveis. \\
      \end{tabularx}}
    \end{column}
    \begin{column}{0.35\textwidth}
      {\footnotesize\color{cinza}\textbf{Parâmetros}}
      \vspace{0.3em}
      {\scriptsize
      \begin{tabularx}{\linewidth}{@{}L{1.0cm}R@{}}
        $o_{it}$        & Carteira (kg). \\
        $\hat{d}_{it}$  & Previsão líquida (kg). \\
        $m_i$           & Margem (R\$/kg). \\
        $v_i$           & Valor do estoque (R\$/kg). \\
        $c_j$           & \textit{Setup} (R\$/h). \\
        $q_i$           & Lote mínimo (kg). \\
        $\theta$        & Carregamento (\%\,a.a.). \\
      \end{tabularx}}
    \end{column}
    \begin{column}{0.31\textwidth}
      {\footnotesize\color{cinza}\textbf{Variáveis}}
      \vspace{0.3em}
      {\scriptsize
      \begin{tabularx}{\linewidth}{@{}L{1.0cm}R@{}}
        $x_{ijt}$ & Quantidade (kg). \\
        $y_{ijt}$ & \textit{Setup}. \\
        $I_{it}$  & Estoque. \\
        $a_{it}$  & Atraso da carteira. \\
        $l_{it}$  & Venda perdida. \\
      \end{tabularx}}

      \vspace{0.8em}
      \begin{center}
        \begin{tikzpicture}[font=\scriptsize,
            b/.style={rectangle,rounded corners=2pt,minimum width=2.9cm,minimum height=0.44cm}]
          \node[b,fill=laranjaEsc!25,draw=laranjaEsc] at (0,0.9) {carteira $o_{it}$};
          \node[b,fill=laranja!20,draw=laranja]       at (0,0.35) {previsão $\hat{d}_{it}$};
          \node[font=\scriptsize,text=cinzaEsc,anchor=north] at (0,-0.62)
            {carteira $+$ previsão $=$ demanda total};
        \end{tikzpicture}
      \end{center}
    \end{column}
  \end{columns}

  \note{A notação muda em três pontos frente ao modelo anterior. O bucket é semana, e o índice de dia
  desaparece deste nível: a ordem interna é decisão do operacional.

  Existe agora variável contínua de quantidade, x i j t. A produção deixa de ser o dia inteiro ou
  nada, e passa a ser quanto se produz, com o lote mínimo q i impondo o piso econômico da corrida.

  E a demanda deixa de ser uma corrente só. Carteira é o que já está vendido, com prazo; previsão
  líquida é o que se espera vender além dela. As duas competem pela mesma capacidade, e a estrutura
  do balanço é que decide quem cede.

  Os parâmetros de custo também se separam: margem para as faltas, valor unitário para o estoque,
  custo horário para o setup. Antes um único custo unitário fazia os três papéis.}
\end{frame}

% ---------------------------------------------------------------------------
\begin{frame}{Todas as parcelas do objetivo estão em reais}

  {\footnotesize
  \begin{equation*}
    Z^{S} = \underbrace{\sum_{i,t} m_i\, a_{it}}_{\text{\footnotesize\color{laranjaEsc}atraso carteira}}
      + \underbrace{\sum_{i,t} m_i\, l_{it}}_{\text{\footnotesize\color{laranja}venda perdida}}
      + \underbrace{\sum_{i,j,t} c_j\, \tau_i\, y_{ijt}}_{\text{\footnotesize\color{cinza}\textit{setup}}}
      + \underbrace{\sum_{i,t} \tfrac{\theta}{52}\, v_i\, \bar{I}_{it}}_{\text{\footnotesize\color{cinza}estoque}}
  \end{equation*}}

  \restricoes{%
    (A) balanço de estoque \hfill (B) horas da máquina \hfill
    (C) lote mínimo e ligação}

  {\footnotesize
  \begin{align*}
    I_{i,t-1} + \sum_{j \in J_i} x_{ijt}
      &= I_{it} + \bigl(o_{it} + a_{i,t-1} - a_{it}\bigr) + \bigl(\hat{d}_{it} - l_{it}\bigr) \tag{A} \\
    \sum_{i \in I} \Bigl( \tfrac{x_{ijt}}{p_{ij}} + \tau_i\, y_{ijt} \Bigr) &\le \kappa_{jt} \tag{B} \\
    q_i\, y_{ijt} \le x_{ijt} &\le M_{ijt}\, y_{ijt} \tag{C}
  \end{align*}}

  \notacao{%
    \sy{i}{balão} · \sy{j}{máquina} · \sy{t}{semana} · \sy{J_i}{máquinas compatíveis} · 
    \sy{x}{produção} · \sy{y}{\textit{setup}} · \sy{I}{estoque} · \sy{\bar{I}}{estoque médio} · 
    \sy{a}{atraso} · \sy{l}{venda perdida} · \sy{o}{carteira} · 
    \sy{\hat{d}}{previsão líquida} · \sy{\kappa}{horas} · 
    \sy{p}{taxa} · \sy{\tau}{horas de \textit{setup}} · \sy{q}{lote mínimo} · \sy{m}{margem} · 
    \sy{v}{valor} · \sy{c}{custo horário} · \sy{\theta}{carregamento}}

  \note{Quatro parcelas, todas em reais, e nenhum peso a calibrar entre elas.

  As duas primeiras são as faltas, ambas cobradas à margem. Não há multiplicador separando uma da
  outra, e mesmo assim a carteira é estritamente mais cara: o atraso entra no somatório a cada
  semana em que persiste, e o quilograma atrasado continua no balanço, ainda terá de ser produzido.
  A venda perdida é cobrada uma vez e sai do balanço, poupando a produção. A hierarquia sai da
  estrutura, não de um peso arbitrado por quem opera o modelo.

  A terceira parcela é o setup como custo direto: horas de troca vezes custo horário da máquina.
  O custo de oportunidade do setup deixou de ser cobrado, como fazia o modelo anterior, porque
  essas horas já consomem capacidade em C, e a produção perdida já reaparece como falta penalizada.
  Cobrar nos dois lugares contava a mesma perda duas vezes.

  A quarta é o carregamento de estoque: taxa anual sobre o valor unitário, rateada por semana,
  aplicada ao estoque médio. É o que dá fundamento financeiro ao trade-off entre setup e estoque.

  Em A o termo da carteira é o que venceu, mais o que vinha atrasado, menos o que segue atrasado:
  a falta de pedido firme não some, é transportada. C é o lote mínimo.

  Não há camada de cobertura. Uma meta semana a semana seria redundante com o horizonte: o modelo
  já enxerga a demanda futura dentro de T e antecipa produção sozinho quando a capacidade aperta.
  O que ela evitaria é só o efeito de borda na última semana, e isso o horizonte rolante resolve.}
\end{frame}

% ---------------------------------------------------------------------------
\begin{frame}{O semanal entrega ao operacional uma meta de estoque}

  \vspace{0.6em}
  \begin{center}
    \resizebox{0.92\textwidth}{!}{%
    \begin{tikzpicture}[font=\small,
        cx/.style={rectangle,rounded corners=3pt,very thick,minimum width=5.4cm,
                   minimum height=1.5cm,text width=5.0cm,align=center,inner sep=6pt},
        ar/.style={draw=cinza,-{Stealth[length=3mm]},very thick}]

      \node[cx,draw=laranja,fill=laranja!18] (sem) at (0,2.2)
        {\textbf{Semanal}\\[2pt]{\footnotesize 13 semanas · balão · máquina}};
      \node[cx,draw=laranjaEsc,fill=laranjaEsc!18] (ope) at (0,-1.6)
        {\textbf{Operacional}\\[2pt]{\footnotesize 1 semana · turno · posição · cor}};

      \path[ar] (sem.south) -- node[right=4pt,font=\footnotesize,text width=4.6cm,align=left]
        {$\Sigma_i$: estoque do balão $i$ ao fim da semana\\[2pt]
         \textcolor{cinza}{restrição \emph{suave}, com folga $g_i$ penalizada}} (ope.north);

      \node[anchor=east,font=\footnotesize,text=cinzaEsc,text width=4.2cm,align=right] at (-3.1,0.3)
        {\textbf{não} desce:\\ máquina, dia, sequência};

      \draw[draw=cinza!60,-{Stealth[length=3mm]},thick,dashed]
        (ope.east) .. controls (6.2,-1.6) and (6.2,2.2) .. (sem.east);
      \node[anchor=west,font=\footnotesize,text=cinza,text width=3.4cm,align=left] at (6.4,0.3)
        {sobe o estado: estoque, carteira consumida, forma e cor montadas};
    \end{tikzpicture}}
  \end{center}

  \note{Este é o contrato entre os dois níveis, e é deliberadamente estreito.

  O que desce é uma única grandeza: quanto de estoque cada balão deve ter ao fim da semana. E
  desce como restrição suave — se a meta for inatingível, o operacional paga a folga e continua
  viável, em vez de devolver inviabilidade.

  O que explicitamente não desce é a alocação. O modelo anterior fixava janelas: bloco de dias em
  que a máquina roda um balão, herdado do tático e imutável. Isso garantia consistência, mas
  tirava do operacional qualquer capacidade de reação: um prazo que não coubesse na janela virava
  atraso, sem alternativa. Agora o operacional escolhe máquina, turno e sequência livremente, e a
  consistência é cobrada pelo estado final, não pelo caminho.

  Na direção contrária sobe o realizado, que alimenta a re-execução seguinte em horizonte
  rolante: estoque, carteira consumida e a configuração em que cada máquina ficou.}
\end{frame}

% ---------------------------------------------------------------------------
\begin{frame}{A troca de cor é direcional, com transições proibidas}

  \vspace{0.4em}
  \begin{columns}[c,onlytextwidth]
    \begin{column}{0.60\textwidth}
      \begin{center}
      \scriptsize
      \setlength{\tabcolsep}{5pt}
      \begin{tabular}{@{}l|ccccc@{}}
        \toprule
        \diagbox[width=2.3cm,height=2.2em]{\textbf{de}}{\textbf{para}}
              & branco & amarelo & verde & vermelho & preto \\
        \midrule
        branco   & --     & 0{,}5  & 0{,}5  & 0{,}8  & 1{,}0 \\
        amarelo  & \cellcolor{laranjaEsc!20}$\times$ & --     & 0{,}6  & 0{,}8  & 1{,}0 \\
        verde    & \cellcolor{laranjaEsc!20}$\times$ & 1{,}6  & --     & 1{,}0  & 1{,}2 \\
        vermelho & \cellcolor{laranjaEsc!20}$\times$ & \cellcolor{laranjaEsc!20}$\times$ & 1{,}8 & -- & 1{,}2 \\
        preto    & \cellcolor{laranjaEsc!20}$\times$ & \cellcolor{laranjaEsc!20}$\times$ & 2{,}0 & 1{,}8 & -- \\
        \bottomrule
      \end{tabular}
      \end{center}
      \vspace{0.5em}
      \begin{center}
        {\scriptsize\color{cinza}horas \qquad \textcolor{laranjaEsc}{$\times$} proibida}
      \end{center}
    \end{column}
    \begin{column}{0.37\textwidth}
      \begin{center}
        \begin{tikzpicture}[font=\footnotesize]
          \node[circle,draw=cinza,thick,fill=white,minimum size=1.0cm] (cl) at (0,1.5) {claro};
          \node[circle,draw=cinza,thick,fill=cinza!75,minimum size=1.0cm,text=white] (es) at (0,-0.5) {escuro};
          \draw[-{Stealth[length=3mm]},very thick,cinza] (-0.55,1.1) to[bend right=35] node[left,font=\scriptsize,text=cinzaEsc] {menor tempo} (-0.55,-0.15);
          \draw[-{Stealth[length=3mm]},very thick,laranja] (0.55,-0.15) to[bend right=35] node[right,font=\scriptsize,text=cinzaEsc] {maior tempo} (0.55,1.1);
        \end{tikzpicture}
      \end{center}
    \end{column}
  \end{columns}

  \note{Esta matriz é o dado que carrega a assimetria. Ler linha para coluna: a célula é o tempo de
  limpeza para sair da cor da linha e entrar na cor da coluna.

  O padrão é claro: a metade de cima consome menos tempo, a de baixo consome mais, e a coluna do
  branco está inteiramente bloqueada. Depois de rodar preto, voltar para branco não é uma questão
  de tempo — é inviável sem purgar o tanque inteiro, e isso não é uma operação de setup, é uma
  parada.

  Na implementação isso é uma matriz DE-PARA lida do Excel, com célula vazia significando
  transição proibida. Existe uma alternativa na literatura, que é modelar setup por
  características, com um gama por tipo em vez de uma matriz de tamanho número de cores ao
  quadrado. É mais compacta e unificaria formato e cor num formalismo só, mas não expressa
  proibição naturalmente. A matriz foi adotada por causa disso, e essa escolha está declarada no
  texto.}
\end{frame}

% ---------------------------------------------------------------------------
\begin{frame}{A forma fica fixa no turno e só a cor muda}

  \begin{columns}[T,onlytextwidth]
    \begin{column}{0.30\textwidth}
      {\footnotesize\color{laranja}\textbf{Conjuntos e parâmetros}}
      \vspace{0.3em}
      {\scriptsize
      \begin{tabularx}{\linewidth}{@{}L{1.2cm}R@{}}
        $S$          & Turnos da semana. \\
        $N$          & Posições no turno. \\
        $K_i$        & Cores do balão $i$. \\
        $A$          & Transições permitidas. \\
        $f_{i'i}$    & \textit{Setup} de forma (h). \\
        $e_{kk'}$    & \textit{Setup} de cor (h). \\
        $b_{js}$     & Horas do turno. \\
        $\Sigma_i$   & Meta de estoque. \\
      \end{tabularx}}
    \end{column}
    \begin{column}{0.30\textwidth}
      {\footnotesize\color{laranja}\textbf{Variáveis}}
      \vspace{0.3em}
      {\scriptsize
      \begin{tabularx}{\linewidth}{@{}L{1.2cm}R@{}}
        $w_{ijs}$      & Forma montada. \\
        $\nu_{kjsn}$   & Cor na posição. \\
        $x_{kjsn}$     & Quantidade (kg). \\
        $z_{i'ijs}$    & Troca de forma. \\
        $\zeta_{kk'jsn}$ & Troca de cor. \\
        $\varphi_{kjs}$ & Cor ao fim do turno. \\
        $a_{ks},\,l_{ks}$ & Atraso, perda. \\
        $g_i$          & Folga da meta. \\
      \end{tabularx}}
    \end{column}
    \begin{column}{0.36\textwidth}
      \vspace{0.6em}
      \begin{center}
        \begin{tikzpicture}[font=\scriptsize]
          \node[rectangle,draw=laranja,very thick,fill=laranja!15,
                minimum width=4.0cm,minimum height=0.62cm] at (0,1.3) {turno $s$ — forma $i$ fixa};
          \node[rectangle,draw=cinza,thick,fill=corBranco,minimum width=1.15cm,minimum height=0.6cm] at (-1.3,0.35) {};
          \node[rectangle,draw=cinza,thick,fill=corAmarelo,minimum width=1.15cm,minimum height=0.6cm] at (0,0.35) {};
          \node[rectangle,draw=cinza,thick,fill=corVermelho,minimum width=1.15cm,minimum height=0.6cm] at (1.3,0.35) {};
          \node[font=\scriptsize,text=cinza] at (-1.3,-0.25) {$n=1$};
          \node[font=\scriptsize,text=cinza] at (0,-0.25)    {$n=2$};
          \node[font=\scriptsize,text=cinza] at (1.3,-0.25)  {$n=3$};
          \draw[-{Stealth[length=2mm]},thick,laranjaEsc] (-0.72,0.35) -- (-0.58,0.35);
          \draw[-{Stealth[length=2mm]},thick,laranjaEsc] (0.58,0.35) -- (0.72,0.35);
          \node[font=\scriptsize,text=cinzaEsc] at (0,-0.95) {a ordem das posições \textbf{é} a sequência};
        \end{tikzpicture}
      \end{center}
    \end{column}
  \end{columns}

  \note{O bucket é o turno, e dentro dele há um número pequeno de posições. A ordem das posições é a
  sequência de produção — não existe um modelo de sequenciamento separado depois.

  A restrição estrutural é que a forma fica fixa no turno: dentro dele só a cor muda. Isso vem da
  escala dos tempos. Uma troca de forma leva de três a sete horas num turno de oito; fazê-la no
  meio do turno não é operação real. Ela acontece na virada, e consome horas do turno que começa.

  É a estrutura de macroperíodo e microperíodos do GLSP, com a família fixada no bloco — o
  block planning de Günther. A economia de binárias vem daí: em vez de sequenciar livremente
  sobre balão e cor, sequencia-se cor dentro de uma forma já decidida.

  A variável phi é a que carrega o estado de cor entre turnos, inclusive através de turnos
  ociosos. Sem ela, uma máquina parada por um turno esqueceria a configuração e atravessaria a
  troca de graça.}
\end{frame}

% ---------------------------------------------------------------------------
\begin{frame}{Seis restrições fecham o modelo por turno}

  \restricoes{%
    (A) uma cor por posição ativa \hfill (B) a cor pertence à forma montada \hfill
    (C) transição proibida \\[0.2em]
    (D) cor persiste entre turnos \hfill (E) horas do turno \hfill
    (F) meta recebida do semanal}

  {\footnotesize
  \begin{align*}
    \sum_{k} \nu_{kjsn} &= u_{js} \tag{A} \\
    \nu_{kjsn} &\le w_{i(k)js} \tag{B} \\
    \nu_{kjs,n-1} + \nu_{k'jsn} &\le 1 \qquad {\scriptstyle (k,k') \notin A} \tag{C} \\
    \bigl| \varphi_{kjs} - \varphi_{kj,s-1} \bigr| &\le u_{js} \tag{D} \\
    \sum_{k,n} \tfrac{x_{kjsn}}{p_{i(k)j}}
      + \sum_{i',i} f_{i'i} z_{i'ijs}
      + \sum_{(k,k') \in A,\, n} e_{kk'} \zeta_{kk'jsn} &\le b_{js} \tag{E} \\
    \sum_{k \in K_i} I_{k,\bar{s}} + g_i &\ge \Sigma_i \tag{F}
  \end{align*}}

  \notacao{%
    \sy{i}{balão} · \sy{j}{máquina} · \sy{s}{turno} · \sy{n}{posição} · \sy{k}{cor} · 
    \sy{K_i}{cores do balão} · \sy{A}{transições permitidas} · \sy{\bar{s}}{último turno} · 
    \sy{\nu}{cor na posição} · \sy{w}{forma montada} · \sy{u}{turno ativo} · \sy{x}{produção} · 
    \sy{z}{troca de forma} · \sy{\zeta}{troca de cor} · \sy{\varphi}{cor ao fim do turno} · 
    \sy{I}{estoque} · \sy{g}{folga} · \sy{f}{\textit{setup} de forma (h)} · 
    \sy{e}{\textit{setup} de cor (h)} · \sy{b}{horas do turno} · \sy{p}{taxa} · 
    \sy{\Sigma}{meta do semanal}}

  \note{Seis restrições dão a estrutura toda.

  A e B dizem que toda posição de um turno ativo recebe exatamente uma cor, e que a cor tem de
  pertencer à forma montada. É o que fixa a família dentro do bloco.

  C é a proibição. A matriz de troca tem células vazias, e onde não há tempo de troca não há
  variável de transição: entra esta desigualdade, que impede as duas cores de se sucederem.

  D é o carryover de estado. Quando a máquina não produz no turno, u é zero e a igualdade força a
  cor a permanecer a do turno anterior. Sem isso, um turno ocioso apaga a configuração e o modelo
  atravessa trocas de graça.

  E é a restrição central: produção mais setups cabem nas horas do turno. É dura, não é relatório
  de excedente. Uma ordem não atravessa a virada.

  F é o contrato com o nível semanal: o estoque somado das cores do balão tem de alcançar a meta,
  com folga penalizada.}
\end{frame}

% ===========================================================================
\section{Produto tecnológico}
% ===========================================================================

\begin{frame}{Os dois modelos rodam num sistema completo e operável}

  \begin{center}
    \resizebox{\textwidth}{!}{%
    \begin{tikzpicture}[font=\small,
        mod/.style={rectangle,rounded corners=2pt,draw=cinza,thick,fill=white,
                    text width=3.3cm,minimum height=1.2cm,text centered,inner sep=4pt,font=\footnotesize},
        ar/.style={draw=cinza,-{Stealth[length=2.5mm]},thick}]

      \node[mod,fill=cinza!12,draw=cinza] (xls) {\textbf{Excel}\\{\scriptsize pedidos, demanda, matrizes}};
      \node[mod,right=1.1cm of xls] (data) {\texttt{data.py}};
      \node[mod,right=1.1cm of data,fill=laranja!15,draw=laranja] (sem) {\texttt{weekly\_model.py}\\{\scriptsize MILP semanal}};
      \node[mod,right=1.1cm of sem,fill=laranjaEsc!15,draw=laranjaEsc] (dia) {\texttt{daily\_model.py}\\{\scriptsize MILP por turno}};

      \node[mod,below=1.4cm of sem] (planner) {\texttt{planner.py}};
      \node[mod,left=1.1cm of planner,fill=cinza!15] (web) {\textbf{FastAPI + React}};
      \node[mod,right=1.1cm of planner,fill=cinza!12] (hist) {\texttt{history.py}\\{\scriptsize cenários salvos}};

      \path[ar] (xls) -- (data);
      \path[ar] (data) -- (sem);
      \path[ar] (sem) -- node[above,font=\scriptsize,text=laranjaEsc]{$\Sigma_i$} (dia);
      \path[ar] (sem.south) -- (planner.north);
      \path[ar] (dia.south) to[out=-90,in=0] (planner.east);
      \path[ar] (planner) -- (web);
      \path[ar] (planner.east) to[out=0,in=180] (hist.west);
    \end{tikzpicture}}
  \end{center}

  \note{O pipeline ficou mais curto que o anterior. Saíram três módulos: o que montava as janelas e
  os dois schedulers de cor intercambiáveis. Com o sequenciamento dentro do modelo diário, não há
  mais etapa de pós-processamento e não há mais heurística a escolher por configuração.

  A superfície entre os dois modelos é uma grandeza só, o sigma: a meta de estoque por balão ao
  fim da semana congelada. O modelo diário não recebe alocação, não recebe janela, não recebe
  ordem — recebe um estado a alcançar.

  A hierarquia que esses módulos realizam tem três níveis: a previsão entrega demanda e o alfa
  para o semanal, e o semanal entrega a meta de estoque de fim de semana para o operacional. Na
  direção contrária, o operacional devolve estoque realizado, atrasos e a configuração em que
  cada máquina ficou. Essa volta existe entre execuções, no horizonte rolante, mas não dentro de
  uma execução: o semanal não revê seu plano à luz do que o operacional descobriu na mesma rodada.

  Na ponta há um dashboard React com backend FastAPI: parametriza o cenário, dispara as duas
  etapas em sequência e grava cada execução no histórico, o que permite comparar cenários depois
  sem rodar de novo. As visualizações são um Gantt de máquina por turno, a natureza da produção —
  contra pedido, contra previsão, contra estoque — e o atendimento da carteira.}
\end{frame}

% ===========================================================================
\section{Conclusão}
% ===========================================================================

\begin{frame}{A decomposição vem do processo, não da conveniência}

  \vspace{1.0em}
  \begin{center}
    \begin{tikzpicture}[font=\small,
        c/.style={rectangle,rounded corners=3pt,thick,draw=cinza,fill=white,text width=5.4cm,
                  minimum height=1.15cm,text centered,inner sep=5pt}]

      \node[c,draw=laranja,very thick,fill=laranja!12] at (-3.1, 1.9) {decomposição vinda do processo};
      \node[c] at ( 3.1, 1.9) {previsão própria integrada};
      \node[c] at (-3.1, 0.4) {\textit{setup} direcional com proibições};
      \node[c] at ( 3.1, 0.4) {custo de oportunidade puro};
      \node[c] at ( 0.0,-1.1) {sistema completo, operável};

    \end{tikzpicture}
  \end{center}

  \note{Cinco pontos.

  Primeiro, a decomposição hierárquica vem do processo físico, não da conveniência computacional.
  A fronteira família-item coincide com a fronteira entre os dois setups. Esse é o argumento
  central.

  Segundo, a previsão de demanda é modelada dentro do trabalho e conectada ao otimizador pelo
  alfa. Toda a literatura de lot sizing revisada recebe a previsão pronta de fora.

  Terceiro, setup direcional com transições proibidas, que nenhuma das quatro classes revisadas
  trata.

  Quarto, objetivo de custo de oportunidade puro, sem holding cost, coerente com a forma como a
  empresa avalia a decisão — ela não capitaliza estoque, ela pensa em máquina parada e venda
  perdida.

  Quinto, e este é o de produto: o trabalho não para na formulação. São dois solvers
  intercambiáveis, um gerador de instâncias e uma interface, tudo rodando.}
\end{frame}

% ---------------------------------------------------------------------------
\begin{frame}{A prioridade deve vir da margem ponderada por cliente}

  \vspace{2.4em}
  \begin{center}
    \begin{tikzpicture}[font=\small,
        cx/.style={rectangle,rounded corners=2pt,draw=cinza,thick,fill=white,inner sep=8pt},
        cn/.style={rectangle,rounded corners=2pt,draw=laranja,very thick,fill=laranja!12,inner sep=8pt}]

      \node[cx] (hoje) at (0,0)
        {$\textstyle\sum_{i,t} m_i \bigl( a_{it} + l_{it} \bigr)$};
      \node[font=\footnotesize,text=cinza,above=3pt of hoje] {hoje};

      \node[cn,right=1.0cm of hoje] (prop)
        {$\textstyle\sum_{i,t} \alert{\omega_{ic}}\, m_i \bigl( a_{it} + l_{it} \bigr)$};
      \node[font=\footnotesize,text=laranjaEsc,above=3pt of prop] {proposto};

      \draw[-{Stealth[length=3mm]},thick,cinza] (hoje) -- (prop);

      \node[font=\scriptsize,text=cinza,anchor=north,text width=4.4cm,align=center]
        at ([yshift=-7pt]hoje.south) {custo unitário único \\ no dado da empresa};
      \node[font=\scriptsize,text=cinzaEsc,anchor=north,text width=5.4cm,align=center]
        at ([yshift=-7pt]prop.south) {margem efetiva \\ peso do gestor por cliente e balão};
    \end{tikzpicture}
  \end{center}

  \note{As duas faltas do objetivo semanal — atraso de carteira e venda perdida
  — são hoje ponderadas por uma única grandeza, o $m$, calculado como preço de venda menos custo
  unitário sobre a planilha que a empresa forneceu. Nessa planilha o custo unitário é o mesmo para
  os dezessete balões, de modo que a diferenciação entre produtos vem só do preço. A prioridade
  que o modelo aplica hoje é, portanto, herdada de um dado de custo que não distingue produto.

  A primeira melhoria é usar margem efetiva por produto, com custo próprio de cada balão. A
  diferença não é de precisão apenas: penalizar a falta pela margem transforma o objetivo em custo
  de oportunidade, e minimizar oportunidade perdida é um proxy para maximizar lucro — enquanto
  minimizar custo apenas reduz gasto, sem olhar o que cada quilo não vendido deixa de trazer.

  A segunda melhoria é o peso destacado. Margem é a dimensão quantitativa, mas a decisão real tem
  uma dimensão qualitativa que ela não captura: um cliente estratégico, um produto que sustenta
  contrato, um lançamento que não pode faltar na prateleira. O ômega é um peso por cliente e
  balão, arbitrado pelo tomador de decisão, que multiplica a penalidade das faltas daquele par.
  Com ele o gestor enviesa deliberadamente o modelo para proteger o nível de serviço do que
  considera importante, sem reescrever formulação e sem falsear a margem.

  O dado necessário já existe: a carteira lida pelo sistema traz o cliente de cada pedido, mas
  esse campo não entra hoje em nenhuma das duas formulações. Incorporá-lo exige indexar a carteira
  e o atraso também por cliente, o que aumenta o modelo semanal.

  Outros pontos seguem em aberto e convém declará-los. A perecibilidade do látex está descrita no
  texto e ausente do modelo: ou se modelam tanque e validade, ou a menção é rebaixada a premissa.
  E não há realimentação do turno para a semana dentro de uma execução — ela ocorre na re-execução
  seguinte, em horizonte rolante.}
\end{frame}

% ---------------------------------------------------------------------------
\begin{frame}[plain]
  \vfill
  \begin{center}
    {\Huge Obrigado}
  \end{center}
  \vfill
\end{frame}

% ---------------------------------------------------------------------------
\begin{frame}[allowframebreaks]{Referências}
  \scriptsize
  \bibliography{referencias}
\end{frame}

\end{document}
